% =====================================================================
% Anexo D: analisis arquitectonico y computacional de los componentes.
% =====================================================================

\section*{Anexo C. Arquitectura del \emph{encoder} de Sonata}
\phantomsection\label{sec:analisis-componentes}\label{sec:arq-sonata}
\addcontentsline{toc}{section}{Anexo C. Arquitectura del \emph{encoder} de Sonata}

% Numeracion propia del anexo: figuras, tablas y ecuaciones como D.n
\setcounter{figure}{0}\renewcommand{\thefigure}{C.\arabic{figure}}
\setcounter{table}{0}\renewcommand{\thetable}{C.\arabic{table}}
\setcounter{equation}{0}\renewcommand{\theequation}{C.\arabic{equation}}

Sonata usa un Point Transformer~V3 como codificador de cinco etapas, sin
decodificador y con $108{,}46$~M de parámetros. Cada bloque combina posición,
atención y capas lineales mediante conexiones residuales; entre etapas, el
\emph{pooling} agrupa puntos de una misma celda, duplica su tamaño y aumenta los
canales para incorporar más contexto. La Figura~\ref{fig:arq-sonata} resume esta
jerarquía y la Tabla~\ref{tab:params-sonata} verifica el total: el 94{,}3\,\% se
concentra en las etapas 4 y 5. Los pesos ocupan $0{,}43$~GB, mientras los
34{,}60\,GB medidos en PL13 provienen principalmente de las activaciones
(Sección~\ref{sec:res-rama-ssl}).

\begin{figure}[!htbp]
\centering
\begin{tikzpicture}[
  font=\footnotesize,
  lbl/.style={draw=black!65, line width=0.4pt, text width=5.1cm, align=center,
              inner sep=3pt, minimum height=1.05cm},
  txt/.style={draw=black!65, line width=0.4pt, text width=8.3cm, align=left,
              inner sep=3pt, minimum height=1.05cm, anchor=west,
              execute at begin node={\hyphenpenalty=10000\exhyphenpenalty=10000 }},
  fl/.style={-{Latex[length=1.5mm]}, line width=0.4pt, black!70},
  tap/.style={line width=0.4pt, black!70}
]
\def\paso{1.4}
\node[lbl, fill=black!5] (l0) at (0,0)
  {\textbf{Entrada}\\[-1pt]{\scriptsize $N$ puntos $\times$ 9 canales}};
\node[txt, fill=black!5] (t0) at (2.65,0)
  {Una curva Z o Hilbert los ordena para que los vecinos queden juntos, y pasan de 9 a 48 canales.};

\node[lbl, fill=blue!5] (l1) at (0,-1*\paso)
  {\textbf{Etapa 1}\\[-1pt]{\scriptsize 3 bloques $\cdot$ 48 canales $\cdot$ 3 cabezas}};
\node[txt, fill=blue!5] (t1) at (2.65,-1*\paso)
  {Detalle fino. Cada punto se relaciona con sus vecinos más próximos, en celdas de 2\,cm.};

\node[lbl, fill=blue!9] (l2) at (0,-2*\paso)
  {\textbf{Etapa 2}\\[-1pt]{\scriptsize 3 bloques $\cdot$ 96 canales $\cdot$ 6 cabezas}};
\node[txt, fill=blue!9] (t2) at (2.65,-2*\paso)
  {Con \emph{pooling} junta los puntos en celdas de 4\,cm y sube a 96 canales.};

\node[lbl, fill=blue!13] (l3) at (0,-3*\paso)
  {\textbf{Etapa 3}\\[-1pt]{\scriptsize 3 bloques $\cdot$ 192 canales $\cdot$ 12 cabezas}};
\node[txt, fill=blue!13] (t3) at (2.65,-3*\paso)
  {Celdas de 8\,cm y 192 canales. Desde aquí las etapas aportan al descriptor final.};

\node[lbl, fill=blue!17] (l4) at (0,-4*\paso)
  {\textbf{Etapa 4}\\[-1pt]{\scriptsize 12 bloques $\cdot$ 384 canales $\cdot$ 24 cabezas}};
\node[txt, fill=blue!17] (t4) at (2.65,-4*\paso)
  {Celdas de 16\,cm y 384 canales. Es la etapa más profunda, con 12 bloques.};

\node[lbl, fill=blue!21] (l5) at (0,-5*\paso)
  {\textbf{Etapa 5}\\[-1pt]{\scriptsize 3 bloques $\cdot$ 512 canales $\cdot$ 32 cabezas}};
\node[txt, fill=blue!21] (t5) at (2.65,-5*\paso)
  {Celdas de 32\,cm y 512 canales. Es la de contexto más amplio.};

\node[lbl, fill=orange!15] (lf) at (0,-6.35*\paso)
  {$\boldsymbol{f_{\mathrm{ssl}}}$\\[-1pt]{\scriptsize $192+384+512=1088$}};
\node[txt, fill=orange!15] (tf) at (2.65,-6.35*\paso)
  {Cada punto recibe lo de las etapas 3, 4 y 5 (\emph{up-cast}), con detalle y contexto a la vez.};

\foreach \a/\b in {0/1,1/2,2/3,3/4,4/5}
  \draw[fl] (l\a.south) -- (l\b.north);

\coordinate (bus) at ($(t3.east)+(0.5,0)$);
\foreach \e in {t3, t4, t5}
  \draw[tap] (\e.east) -- (\e.east -| bus);
\draw[tap] (t3.east -| bus) -- (tf.east -| bus);
\draw[fl] (tf.east -| bus) -- (tf.east);
\end{tikzpicture}
\caption[Arquitectura del \emph{encoder} de Sonata]{Arquitectura del \emph{encoder} de Sonata, cinco etapas con celdas cada vez mayores. Las etapas 3, 4 y 5 forman $f_{\mathrm{ssl}}$.}
\label{fig:arq-sonata}
\end{figure}

\begin{table}[!htbp]
\centering
\setstretch{1}\small
\caption[Parámetros del \emph{encoder} de Sonata por etapa]{Distribución de parámetros del \emph{encoder} de Sonata por etapa.}
\label{tab:params-sonata}
\begin{tabular}{@{}lrrrr@{}}
\toprule
\textbf{Etapa} & \textbf{Canales} & \textbf{Bloques} & \textbf{Parámetros} & \textbf{\%} \\
\midrule
1 & 48 & 3 & $278\,928$ & 0{,}3 \\
2 & 96 & 3 & $1\,110\,816$ & 1{,}0 \\
3 & 192 & 3 & $4\,433\,472$ & 4{,}1 \\
4 & 384 & 12 & $70\,857\,216$ & \textbf{65{,}3} \\
5 & 512 & 3 & $31\,483\,392$ & 29{,}0 \\
\emph{Embedding} y submuestreo & & & $297\,504$ & 0{,}3 \\
\midrule
\textbf{Total} & & & $\mathbf{108\,461\,328}$ & 100{,}0 \\
\bottomrule
\end{tabular}
\end{table}
\FloatBarrier
